\documentclass[12pt,a4paper]{article}
\usepackage[utf8]{inputenc}
\usepackage{geometry}
\geometry{a4paper, margin=1in}
\usepackage{hyperref}
\hypersetup{colorlinks=true, linkcolor=blue, urlcolor=blue, citecolor=blue}
\usepackage{graphicx}
\usepackage{listings}
\usepackage{xcolor}
\usepackage{booktabs}
\usepackage{tabularx}
\usepackage{longtable}
\usepackage{array}
\usepackage{tikz}
\usepackage{mdframed}
\usepackage{enumitem}
\usepackage{fancyhdr}
\usepackage{titlesec}
\usetikzlibrary{shapes.geometric, arrows, positioning}

% ── Colours ────────────────────────────────────────────────────────────────
\definecolor{codegreen}{rgb}{0,0.6,0}
\definecolor{codegray}{rgb}{0.5,0.5,0.5}
\definecolor{codepurple}{rgb}{0.58,0,0.82}
\definecolor{backcolour}{rgb}{0.95,0.95,0.92}
\definecolor{passgreen}{rgb}{0.13,0.55,0.13}
\definecolor{failred}{rgb}{0.8,0.1,0.1}
\definecolor{infobox}{rgb}{0.93,0.96,1.0}
\definecolor{warnbox}{rgb}{1.0,0.97,0.87}
\definecolor{headerblue}{rgb}{0.18,0.31,0.60}

% ── Code listing style ─────────────────────────────────────────────────────
\lstdefinestyle{javaStyle}{
    backgroundcolor=\color{backcolour},
    commentstyle=\color{codegreen},
    keywordstyle=\color{magenta}\bfseries,
    numberstyle=\tiny\color{codegray},
    stringstyle=\color{codepurple},
    basicstyle=\ttfamily\scriptsize,
    breakatwhitespace=false,
    breaklines=true,
    captionpos=b,
    keepspaces=true,
    numbers=left,
    numbersep=5pt,
    showspaces=false,
    showstringspaces=false,
    showtabs=false,
    tabsize=2,
    frame=single,
    rulecolor=\color{codegray}
}
\lstdefinestyle{xmlStyle}{
    backgroundcolor=\color{backcolour},
    basicstyle=\ttfamily\scriptsize,
    keywordstyle=\color{blue},
    stringstyle=\color{codepurple},
    commentstyle=\color{codegreen},
    breaklines=true,
    frame=single,
    numbers=left,
    numberstyle=\tiny\color{codegray},
    numbersep=5pt,
    tabsize=2
}
\lstdefinestyle{shellStyle}{
    backgroundcolor=\color{black!90},
    basicstyle=\ttfamily\scriptsize\color{white},
    breaklines=true,
    frame=single,
    rulecolor=\color{black}
}
\lstdefinestyle{gherkinStyle}{
    backgroundcolor=\color{backcolour},
    basicstyle=\ttfamily\scriptsize,
    keywordstyle=\color{blue}\bfseries,
    stringstyle=\color{codepurple},
    commentstyle=\color{codegreen},
    breaklines=true,
    frame=single,
    numbers=left,
    numberstyle=\tiny\color{codegray},
    numbersep=5pt
}
\lstset{style=javaStyle}

% ── Info/warn boxes ────────────────────────────────────────────────────────
\newmdenv[backgroundcolor=infobox, linecolor=headerblue, linewidth=1.5pt,
  innerleftmargin=10pt, innerrightmargin=10pt, innertopmargin=6pt,
  innerbottommargin=6pt]{infoblock}
\newmdenv[backgroundcolor=warnbox, linecolor=orange!70, linewidth=1.5pt,
  innerleftmargin=10pt, innerrightmargin=10pt, innertopmargin=6pt,
  innerbottommargin=6pt]{warnblock}

% ── Page style ─────────────────────────────────────────────────────────────
\pagestyle{fancy}
\fancyhf{}
\fancyhead[L]{\small\textcolor{headerblue}{\textbf{SE3010 -- Student Management System}}}
\fancyhead[R]{\small\textcolor{codegray}{Software Quality \& Testing}}
\fancyfoot[C]{\thepage}
\renewcommand{\headrulewidth}{0.4pt}

\title{
  \vspace{-1cm}
  \textcolor{headerblue}{\Huge\textbf{Student Management System}}\\[0.4cm]
  \Large Testing \& Operations Manual\\[0.3cm]
  \normalsize SE3010 -- Software Quality \& Project Management\\[0.2cm]
  \normalsize Group Assignment Documentation
}
\author{Group Assignment -- Software Testing (TestNG + Cucumber)}
\date{March 3, 2026}

% ═══════════════════════════════════════════════════════════════════════════
\begin{document}

\maketitle
\thispagestyle{empty}

\begin{infoblock}
\textbf{Quick Navigation:} This document covers the complete project -- architecture,
all test cases per member, step-by-step run instructions, and how to verify results.
Use the Table of Contents below to jump to any section.
\end{infoblock}

\vspace{0.5cm}
\tableofcontents
\newpage

% ═══════════════════════════════════════════════════════════════════════════
\section{Project Overview}
% ═══════════════════════════════════════════════════════════════════════════

This document is the complete reference for the \textbf{Student Management System} --
a full-stack web application used as a demonstration project for enterprise-grade software
testing in a university assignment context.

\subsection{Technology Stack}

\begin{tabularx}{\textwidth}{>{\bfseries}lX}
\toprule
Layer & Technology \\
\midrule
Backend & Java 17, Spring Boot 3.4.x, Spring Web, Spring Data JPA, Lombok \\
Database & H2 In-Memory (tests) / Neon PostgreSQL (production) \\
Frontend & React 19, Vite 7, Vanilla CSS \\
Testing Framework & \textbf{TestNG 7.10.2} (all unit \& integration tests) \\
BDD & \textbf{Cucumber 7.18.0} with \texttt{cucumber-testng} runner \\
Mocking & \textbf{Mockito 5.12.0} standalone (no JUnit binding) \\
API Testing & Spring MockMvc (no Selenium, no Postman, no browser) \\
Build / Reports & Maven 3.x, maven-surefire-plugin 3.2.5, maven-surefire-report-plugin \\
\bottomrule
\end{tabularx}

\vspace{0.4cm}
\begin{warnblock}
\textbf{Important:} This project uses \textbf{TestNG only}. JUnit Jupiter,
JUnit Vintage, and \texttt{junit-platform-suite} are all \textbf{explicitly excluded}
from the POM. No \texttt{org.junit.*} imports exist anywhere in the test source.
\end{warnblock}

\subsection{Project Folder Structure}

\begin{lstlisting}[style=shellStyle, numbers=none]
SEPQM/
  frontend/                        <- React + Vite UI
    src/
      App.jsx                      <- Main component (CRUD UI)
      api.js                       <- Fetch API helpers
      index.css                    <- Full custom stylesheet
  se3010-testing/                  <- Spring Boot backend + tests
    src/
      main/java/.../
        StudentController.java     <- REST endpoints
        StudentService.java        <- Business logic + duplicate check
        StudentRepository.java     <- JPA repo (findByEmail added)
        Student.java               <- JPA entity
      test/java/
        tests/
          StudentControllerTest.java  <- Member 1: TestNG Assertions
          StudentServiceTest.java     <- Member 2: Fixtures + DataProvider
          CucumberRunnerTest.java     <- Member 4: BDD Runner
        steps/
          StudentSteps.java           <- Member 3: Cucumber Step Defs
      test/resources/
        features/
          student_registration.feature  <- Member 3: Gherkin scenarios
        application-test.properties     <- H2 in-memory DB config
    testng.xml                       <- TestNG suite (all 4 members)
    pom.xml
\end{lstlisting}

% ═══════════════════════════════════════════════════════════════════════════
\section{System Architecture}
% ═══════════════════════════════════════════════════════════════════════════

\subsection{High-Level Data Flow}

\tikzstyle{iobox}   = [trapezium, trapezium left angle=70, trapezium right angle=110,
  minimum width=3.2cm, minimum height=1cm, text centered, draw=black, fill=blue!25]
\tikzstyle{procbox} = [rectangle, minimum width=3.2cm, minimum height=1cm,
  text centered, draw=black, fill=orange!25]
\tikzstyle{dbbox}   = [rectangle, rounded corners, minimum width=3.2cm,
  minimum height=1cm, text centered, draw=black, fill=red!20]
\tikzstyle{arrow}   = [thick,->,>=stealth]

\begin{figure}[h]
\centering
\begin{tikzpicture}[node distance=2.2cm]
\node (ui)         [iobox]   {React Frontend (Vite)};
\node (controller) [procbox, below of=ui]         {StudentController (REST)};
\node (service)    [procbox, below of=controller] {StudentService (Logic)};
\node (repo)       [procbox, below of=service]    {StudentRepository (JPA)};
\node (db)         [dbbox,   right of=repo, xshift=4.5cm] {PostgreSQL / H2};

\draw [arrow] (ui)         -- node[anchor=west,font=\small] {HTTP GET/POST/PUT/DELETE} (controller);
\draw [arrow] (controller) -- (service);
\draw [arrow] (service)    -- (repo);
\draw [arrow] (repo)       -- node[anchor=north,font=\small] {Hibernate ORM} (db);
\end{tikzpicture}
\caption{System Architecture -- Request Flow}
\end{figure}

\subsection{REST API Endpoints}

\begin{tabularx}{\textwidth}{llX}
\toprule
\textbf{Method} & \textbf{Endpoint} & \textbf{Description} \\
\midrule
POST   & \texttt{/api/students}      & Create new student (409 if email duplicate) \\
GET    & \texttt{/api/students}      & Get all students \\
GET    & \texttt{/api/students/\{id\}} & Get student by ID (404 if not found) \\
PUT    & \texttt{/api/students/\{id\}} & Update student (404 if not found) \\
DELETE & \texttt{/api/students/\{id\}} & Delete student \\
\bottomrule
\end{tabularx}

\subsection{Student Domain Model}

\begin{tabularx}{\textwidth}{llX}
\toprule
\textbf{Field} & \textbf{Type} & \textbf{Description} \\
\midrule
\texttt{id}     & Long   & Auto-generated primary key \\
\texttt{name}   & String & Full name of the student \\
\texttt{email}  & String & Campus email (must be unique) \\
\texttt{age}    & int    & Student age \\
\texttt{course} & String & Enrolled course code (e.g.\ SE3010) \\
\bottomrule
\end{tabularx}

% ═══════════════════════════════════════════════════════════════════════════
\section{Prerequisites \& Environment Setup}
% ═══════════════════════════════════════════════════════════════════════════

\subsection{Required Software}

\begin{tabularx}{\textwidth}{>{\bfseries}lX}
\toprule
Tool & Version / Notes \\
\midrule
Java JDK & 17 or higher (\texttt{java -version} to verify) \\
Apache Maven & 3.8+ (\texttt{mvn -version} to verify) \\
Node.js & 18+ (\texttt{node -v} to verify) -- only for frontend \\
npm & 9+ (bundled with Node.js) \\
Git & Any recent version \\
IDE (optional) & IntelliJ IDEA recommended; import as Maven project \\
\bottomrule
\end{tabularx}

\subsection{Verify Prerequisites}

\begin{lstlisting}[style=shellStyle, numbers=none, caption={Verify all tools are installed}]
# Check Java
java -version
# Expected: openjdk version "17.x.x" or higher

# Check Maven
mvn -version
# Expected: Apache Maven 3.x.x

# Check Node + npm (frontend only)
node -v
npm -v
\end{lstlisting}

% ═══════════════════════════════════════════════════════════════════════════
\section{How to Run -- Backend \& Tests}
% ═══════════════════════════════════════════════════════════════════════════

\subsection{Step 1 -- Navigate to the Backend Project}

\begin{lstlisting}[style=shellStyle, numbers=none]
cd "C:\C Projects\SEPQM\se3010-testing"
\end{lstlisting}

\subsection{Step 2 -- Run the Full Test Suite}

This single command compiles all sources, resolves dependencies, runs every
TestNG and Cucumber test, and generates the HTML report:

\begin{lstlisting}[style=shellStyle, numbers=none, caption={Run all tests}]
mvn test
\end{lstlisting}

\begin{infoblock}
\textbf{What happens:} Maven reads \texttt{testng.xml}, which executes three test groups
in sequence -- Member 1 controller tests, Member 2 service tests, and Member 3+4 BDD tests.
All 13 tests should complete in roughly 30--60 seconds.
\end{infoblock}

\subsection{Step 3 -- Run Only One Member's Tests}

To run a single member's tests in isolation, use the \texttt{-Dtest} flag with the
class name, or run individual TestNG groups:

\begin{lstlisting}[style=shellStyle, numbers=none, caption={Run individual member tests}]
# Member 1 only
mvn test -Dtest=tests.StudentControllerTest

# Member 2 only
mvn test -Dtest=tests.StudentServiceTest

# Member 3+4 only (BDD)
mvn test -Dtest=tests.CucumberRunnerTest
\end{lstlisting}

\subsection{Step 4 -- Run the Spring Boot Application (Production)}

To run the backend server (connects to PostgreSQL, not H2):

\begin{lstlisting}[style=shellStyle, numbers=none, caption={Start the backend server}]
mvn spring-boot:run
# Server starts at: http://localhost:8080
# API base URL:     http://localhost:8080/api/students
\end{lstlisting}

\subsection{Step 5 -- Generate the HTML Surefire Report Separately}

\begin{lstlisting}[style=shellStyle, numbers=none, caption={Generate HTML report only}]
mvn surefire-report:report
# Report location: target/site/surefire-report.html
\end{lstlisting}

% ═══════════════════════════════════════════════════════════════════════════
\section{How to Run -- Frontend}
% ═══════════════════════════════════════════════════════════════════════════

\subsection{Step 1 -- Install Dependencies}

\begin{lstlisting}[style=shellStyle, numbers=none, caption={Install frontend npm packages}]
cd "C:\C Projects\SEPQM\frontend"
npm install
\end{lstlisting}

\subsection{Step 2 -- Start the Development Server}

\begin{lstlisting}[style=shellStyle, numbers=none, caption={Start Vite dev server}]
npm run dev
# Frontend runs at: http://localhost:5173
# Make sure backend is also running at :8080
\end{lstlisting}

\subsection{Step 3 -- Build for Production}

\begin{lstlisting}[style=shellStyle, numbers=none, caption={Build frontend for production}]
npm run build
# Output: frontend/dist/
\end{lstlisting}

\subsection{Step 4 -- Lint Check}

\begin{lstlisting}[style=shellStyle, numbers=none, caption={Run ESLint}]
npm run lint
# Zero errors expected (React Hooks + React Refresh rules enforced)
\end{lstlisting}

\begin{infoblock}
\textbf{Frontend Features:} The React UI supports adding a student (POST), listing all
enrolled students (GET), and deleting a student (DELETE) with a confirmation dialog.
All API calls are in \texttt{src/api.js} and target \texttt{http://localhost:8080/api/students}.
\end{infoblock}

% ═══════════════════════════════════════════════════════════════════════════
\section{Test Cases -- Member 1: TestNG Assertions}
% ═══════════════════════════════════════════════════════════════════════════

\textbf{File:} \texttt{src/test/java/tests/StudentControllerTest.java}\\
\textbf{Approach:} Spring Boot integration test using \texttt{MockMvc} to call the real
\texttt{StudentController} without starting a full HTTP server. Assertions use
\texttt{org.testng.Assert} exclusively.

\subsection{Test Case Summary}

\begin{tabularx}{\textwidth}{>{\ttfamily\small}lX>{\bfseries}l}
\toprule
\textbf{Test Method} & \textbf{What it Verifies} & \textbf{Expected} \\
\midrule
testCreateStudent\_Returns201AndCorrectBody
  & POST valid student payload; asserts HTTP 201,
    JSONPath \texttt{\$.name} = ``Alice'', \texttt{\$.email} = ``alice@test.com'',
    response body not empty
  & \textcolor{passgreen}{PASS} \\
\addlinespace
testGetStudent\_Returns200
  & GET existing student by ID; asserts HTTP 200,
    JSONPath \texttt{\$.name} = ``Bob'', body not null
  & \textcolor{passgreen}{PASS} \\
\addlinespace
testGetStudent\_NotFound\_Returns404
  & GET non-existent ID 99999; asserts HTTP status code equals 404
  & \textcolor{passgreen}{PASS} \\
\addlinespace
testCreateStudent\_MalformedJson\_Returns400
  & POST \texttt{\{\ invalid json\ \}} string; asserts HTTP status code equals 400
  & \textcolor{passgreen}{PASS} \\
\bottomrule
\end{tabularx}

\subsection{Key Annotations Used}

\begin{itemize}[noitemsep]
  \item \texttt{@SpringBootTest} -- loads full application context with MockMvc
  \item \texttt{@AutoConfigureMockMvc} -- auto-wires MockMvc
  \item \texttt{@ActiveProfiles("test")} -- activates H2 in-memory DB
  \item \texttt{@BeforeMethod} -- clears repository before every test
  \item \texttt{@Test} from \texttt{org.testng.annotations}
\end{itemize}

\subsection{Code Listing}

\begin{lstlisting}[language=Java, style=javaStyle, caption={StudentControllerTest.java -- Member 1}]
@SpringBootTest(classes = Se3010TestingApplication.class,
        webEnvironment = SpringBootTest.WebEnvironment.MOCK)
@AutoConfigureMockMvc
@ActiveProfiles("test")
public class StudentControllerTest extends AbstractTestNGSpringContextTests {

    @BeforeMethod
    public void cleanUp() { studentRepository.deleteAll(); }

    @Test
    public void testCreateStudent_Returns201AndCorrectBody() throws Exception {
        Student student = Student.builder()
                .name("Alice").email("alice@test.com").age(22).course("SE3010").build();
        String body = mockMvc.perform(post("/api/students")
                        .contentType(MediaType.APPLICATION_JSON)
                        .content(objectMapper.writeValueAsString(student)))
                .andExpect(status().isCreated())
                .andExpect(jsonPath("$.name").value("Alice"))
                .andReturn().getResponse().getContentAsString();
        Assert.assertFalse(body.isEmpty(), "Response body must not be empty");
    }

    @Test
    public void testGetStudent_NotFound_Returns404() throws Exception {
        int statusCode = mockMvc.perform(get("/api/students/{id}", 99999L))
                .andReturn().getResponse().getStatus();
        Assert.assertEquals(statusCode, 404, "Expected HTTP 404 for missing student");
    }
}
\end{lstlisting}

\subsection{How to Check Member 1 Results}

\begin{lstlisting}[style=shellStyle, numbers=none]
mvn test -Dtest=tests.StudentControllerTest
\end{lstlisting}

\textbf{Expected console output:}

\begin{lstlisting}[style=shellStyle, numbers=none]
[INFO] Tests run: 4, Failures: 0, Errors: 0, Skipped: 0
\end{lstlisting}

\textbf{Report file:} \texttt{target/surefire-reports/tests.StudentControllerTest.txt}

% ═══════════════════════════════════════════════════════════════════════════
\section{Test Cases -- Member 2: TestNG Fixtures \& DataProvider}
% ═══════════════════════════════════════════════════════════════════════════

\textbf{File:} \texttt{src/test/java/tests/StudentServiceTest.java}\\
\textbf{Approach:} Pure unit tests of \texttt{StudentService} using Mockito standalone.
Fresh mocks are created in \texttt{@BeforeMethod} and verified with
\texttt{verifyNoMoreInteractions} in \texttt{@AfterMethod}.
A \texttt{@DataProvider} feeds three student records into the parameterised create test.

\subsection{Test Case Summary}

\begin{tabularx}{\textwidth}{>{\ttfamily\small}lX>{\bfseries}l}
\toprule
\textbf{Test Method} & \textbf{What it Verifies} & \textbf{Expected} \\
\midrule
testCreateStudent [Alice]
  & Stubs \texttt{findByEmail} $\to$ empty, \texttt{save} $\to$ saved object;
    asserts ID not null, name = ``Alice'', email match, verifies both repo calls
  & \textcolor{passgreen}{PASS} \\
\addlinespace
testCreateStudent [Bob]
  & Same as above with Bob's data (DataProvider row 2)
  & \textcolor{passgreen}{PASS} \\
\addlinespace
testCreateStudent [Carol]
  & Same as above with Carol's data (DataProvider row 3)
  & \textcolor{passgreen}{PASS} \\
\addlinespace
testGetStudentById\_Found
  & Stubs \texttt{findById(1L)} $\to$ Optional of Dave;
    asserts present = true, name = ``Dave''
  & \textcolor{passgreen}{PASS} \\
\addlinespace
testGetStudentById\_NotFound
  & Stubs \texttt{findById(99L)} $\to$ Optional.empty();
    asserts present = false
  & \textcolor{passgreen}{PASS} \\
\addlinespace
testGetAllStudents\_ReturnsList
  & Stubs \texttt{findAll} $\to$ list of 1 student;
    asserts list size = 1, name = ``Eve''
  & \textcolor{passgreen}{PASS} \\
\bottomrule
\end{tabularx}

\subsection{Key Annotations Used}

\begin{itemize}[noitemsep]
  \item \texttt{@BeforeMethod} -- creates fresh \texttt{mock(StudentRepository.class)} and wires service
  \item \texttt{@AfterMethod} -- calls \texttt{verifyNoMoreInteractions} to catch unexpected calls
  \item \texttt{@DataProvider(name = "studentData")} -- supplies 3 rows of test data
  \item \texttt{@Test(dataProvider = "studentData")} -- parameterised test execution
\end{itemize}

\subsection{DataProvider Definition}

\begin{lstlisting}[language=Java, style=javaStyle, caption={DataProvider in StudentServiceTest.java}]
@DataProvider(name = "studentData")
public Object[][] studentDataProvider() {
    return new Object[][] {
        {"Alice", "alice@test.com", 22, "SE3010"},
        {"Bob",   "bob@test.com",   25, "SE3020"},
        {"Carol", "carol@test.com", 20, "SE3030"},
    };
}

@Test(dataProvider = "studentData")
public void testCreateStudent(String name, String email, int age, String course) {
    when(studentRepository.findByEmail(email)).thenReturn(Optional.empty());
    when(studentRepository.save(any(Student.class))).thenReturn(saved);

    Student result = studentService.createStudent(input);

    Assert.assertEquals(result.getName(), name, "Name should match");
    verify(studentRepository).findByEmail(email);
    verify(studentRepository).save(any(Student.class));
}
\end{lstlisting}

\subsection{Fixture Lifecycle}

\begin{lstlisting}[language=Java, style=javaStyle, caption={BeforeMethod and AfterMethod fixtures}]
@BeforeMethod
public void setUp() {
    studentRepository = mock(StudentRepository.class);
    studentService    = new StudentService(studentRepository);
    System.out.println("[FIXTURE] Fresh mocks created - test starting");
}

@AfterMethod
public void tearDown() {
    verifyNoMoreInteractions(studentRepository);
    System.out.println("[FIXTURE] Verified - test finished");
}
\end{lstlisting}

\subsection{How to Check Member 2 Results}

\begin{lstlisting}[style=shellStyle, numbers=none]
mvn test -Dtest=tests.StudentServiceTest
\end{lstlisting}

\textbf{Expected console output:}

\begin{lstlisting}[style=shellStyle, numbers=none]
[FIXTURE] Fresh mocks created - test starting
[FIXTURE] Verified - test finished
... (repeated for each of 6 test invocations)
[INFO] Tests run: 6, Failures: 0, Errors: 0, Skipped: 0
\end{lstlisting}

\textbf{Report file:} \texttt{target/surefire-reports/tests.StudentServiceTest.txt}

% ═══════════════════════════════════════════════════════════════════════════
\section{Test Cases -- Member 3: Cucumber BDD (Gherkin Scenarios)}
% ═══════════════════════════════════════════════════════════════════════════

\textbf{Feature file:} \texttt{src/test/resources/features/student\_registration.feature}\\
\textbf{Step definitions:} \texttt{src/test/java/steps/StudentSteps.java}\\
\textbf{Approach:} Business-readable Gherkin scenarios executed via Cucumber
with \texttt{MockMvc} for real API calls against the Spring context (no browser, no Selenium).

\subsection{Gherkin Feature File}

\begin{lstlisting}[style=gherkinStyle, caption={student\_registration.feature -- Member 3}]
Feature: Student Registration API
  As an admin
  I want to register students via the API
  So that I can manage the student roster

  Scenario: Successful student registration
    Given no students exist in the system
    When I register a student with name "Alice" email
         "alice@uni.com" age 22 course "SE3010"
    Then the registration response status should be 201
    And the registered student name should be "Alice"

  Scenario: Duplicate email returns conflict error
    Given a student with email "bob@uni.com" already exists
    When I register a student with name "Bob2" email
         "bob@uni.com" age 24 course "SE3020"
    Then the registration response status should be 409

  Scenario: Missing required field returns bad request error
    Given no students exist in the system
    When I send a registration request with malformed JSON body
    Then the registration response status should be 400
\end{lstlisting}

\subsection{Scenario Summary}

\begin{tabularx}{\textwidth}{lX>{\bfseries}l}
\toprule
\textbf{Scenario} & \textbf{What it Verifies} & \textbf{Expected} \\
\midrule
Successful student registration
  & POST valid student; HTTP 201 returned; response body contains name ``Alice''
  & \textcolor{passgreen}{PASS} \\
\addlinespace
Duplicate email returns conflict error
  & Pre-seed email ``bob@uni.com''; POST same email; HTTP 409 Conflict returned
  & \textcolor{passgreen}{PASS} \\
\addlinespace
Missing required field returns bad request
  & POST \texttt{\{\ bad json\ \}}; Spring rejects with HTTP 400 Bad Request
  & \textcolor{passgreen}{PASS} \\
\bottomrule
\end{tabularx}

\subsection{Step Definition Highlights}

\begin{lstlisting}[language=Java, style=javaStyle, caption={StudentSteps.java -- selected step defs}]
@Given("a student with email {string} already exists")
public void a_student_with_email_already_exists(String email) {
    studentRepository.deleteAll();
    studentRepository.save(
        Student.builder().name("Existing").email(email)
                         .age(20).course("SE0000").build());
}

@When("I register a student with name {string} email {string} age {int} course {string}")
public void i_register_a_student(String name, String email,
                                  int age, String course) throws Exception {
    Student s = Student.builder().name(name).email(email)
                                 .age(age).course(course).build();
    lastResult = mockMvc.perform(post("/api/students")
            .contentType(MediaType.APPLICATION_JSON)
            .content(objectMapper.writeValueAsString(s)))
        .andReturn();
}

@Then("the registration response status should be {int}")
public void the_registration_status_should_be(int expectedStatus) {
    int actual = lastResult.getResponse().getStatus();
    Assert.assertEquals(actual, expectedStatus,
        "Expected HTTP " + expectedStatus + " but got " + actual);
}
\end{lstlisting}

\subsection{How to Check Member 3 Results}

\begin{lstlisting}[style=shellStyle, numbers=none]
mvn test -Dtest=tests.CucumberRunnerTest
\end{lstlisting}

\textbf{Expected console output:}

\begin{lstlisting}[style=shellStyle, numbers=none]
Scenario: Successful student registration         # features/student_registration.feature:7
Scenario: Duplicate email returns conflict error  # features/student_registration.feature:14
Scenario: Missing required field returns bad req  # features/student_registration.feature:20
[INFO] Tests run: 3, Failures: 0, Errors: 0, Skipped: 0
\end{lstlisting}

% ═══════════════════════════════════════════════════════════════════════════
\section{Test Cases -- Member 4: Cucumber TestNG Runner \& Reports}
% ═══════════════════════════════════════════════════════════════════════════

\textbf{File:} \texttt{src/test/java/tests/CucumberRunnerTest.java}\\
\textbf{Approach:} Extends \texttt{AbstractTestNGCucumberTests} (the official Cucumber TestNG
integration class) and configures \texttt{@CucumberOptions} to point at the feature file,
glue package, and output plugins.

\subsection{Runner Configuration}

\begin{lstlisting}[language=Java, style=javaStyle, caption={CucumberRunnerTest.java -- Member 4}]
@CucumberOptions(
    features  = "src/test/resources/features/student_registration.feature",
    glue      = {"steps"},
    plugin    = {
        "pretty",
        "html:target/cucumber-reports/cucumber.html",
        "json:target/cucumber-reports/cucumber.json"
    },
    monochrome = true
)
public class CucumberRunnerTest extends AbstractTestNGCucumberTests {

    @Override
    @DataProvider(parallel = false)
    public Object[][] scenarios() {
        return super.scenarios();
    }
}
\end{lstlisting}

\subsection{TestNG Suite File -- \texttt{testng.xml}}

\begin{lstlisting}[style=xmlStyle, caption={testng.xml -- orchestrates all 4 members}]
<suite name="Student API Test Suite" verbose="2">

  <test name="Member1 - Controller Assertions">
    <classes>
      <class name="tests.StudentControllerTest"/>
    </classes>
  </test>

  <test name="Member2 - Service Fixtures">
    <classes>
      <class name="tests.StudentServiceTest"/>
    </classes>
  </test>

  <test name="Member3+4 - Cucumber BDD">
    <classes>
      <class name="tests.CucumberRunnerTest"/>
    </classes>
  </test>

</suite>
\end{lstlisting}

\subsection{How to Check Member 4 Reports}

\subsubsection{Option A -- Run tests and open the HTML report}

\begin{lstlisting}[style=shellStyle, numbers=none]
# Step 1: Run all tests
mvn test

# Step 2: Open the Surefire HTML report
start target\site\surefire-report.html
\end{lstlisting}

\subsubsection{Option B -- Open the Cucumber-specific HTML report}

\begin{lstlisting}[style=shellStyle, numbers=none]
start target\cucumber-reports\cucumber.html
\end{lstlisting}

\subsubsection{Option C -- Open the TestNG native report}

\begin{lstlisting}[style=shellStyle, numbers=none]
start target\surefire-reports\index.html
\end{lstlisting}

\subsection{Report File Locations Summary}

\begin{tabularx}{\textwidth}{>{\ttfamily\small}lX}
\toprule
\textbf{File Path} & \textbf{Description} \\
\midrule
target/site/surefire-report.html & Maven Surefire HTML report (all members) \\
target/surefire-reports/index.html & TestNG HTML suite results \\
target/surefire-reports/emailable-report.html & TestNG single-page email report \\
target/surefire-reports/testng-results.xml & TestNG raw XML results \\
target/cucumber-reports/cucumber.html & Cucumber BDD HTML report (Member 3+4) \\
target/cucumber-reports/cucumber.json & Cucumber JSON (for CI integrations) \\
target/surefire-reports/tests.StudentControllerTest.txt & Member 1 text log \\
target/surefire-reports/tests.StudentServiceTest.txt & Member 2 text log \\
\bottomrule
\end{tabularx}

% ═══════════════════════════════════════════════════════════════════════════
\section{Complete Test Results Summary}
% ═══════════════════════════════════════════════════════════════════════════

Running \texttt{mvn test} from the project root executes all 13 test invocations
across 4 members. The expected final output is:

\begin{lstlisting}[style=shellStyle, numbers=none, caption={Expected final mvn test output}]
[INFO] Tests run: 13, Failures: 0, Errors: 0, Skipped: 0
[INFO] BUILD SUCCESS
\end{lstlisting}

\begin{tabularx}{\textwidth}{lllX}
\toprule
\textbf{Member} & \textbf{Class} & \textbf{Tests} & \textbf{Description} \\
\midrule
Member 1 & StudentControllerTest & 4 & HTTP 201, 200, 404, 400 via MockMvc \\
Member 2 & StudentServiceTest    & 6 & 3$\times$DataProvider + 3 service unit tests \\
Member 3+4 & CucumberRunnerTest  & 3 & BDD: 201, 409, 400 Cucumber scenarios \\
\midrule
\multicolumn{2}{l}{\textbf{Total}} & \textbf{13} & \textbf{All pass, 0 failures} \\
\bottomrule
\end{tabularx}

% ═══════════════════════════════════════════════════════════════════════════
\section{Frontend Application}
% ═══════════════════════════════════════════════════════════════════════════

\subsection{Overview}

The React frontend (\texttt{frontend/}) provides a single-page dashboard to manage student
enrolments. It communicates with the Spring Boot backend via the browser Fetch API.
ESLint (React Hooks + React Refresh rules) reports \textbf{zero warnings or errors}.

\subsection{Component Structure}

\begin{tabularx}{\textwidth}{>{\ttfamily\small}lX}
\toprule
\textbf{File} & \textbf{Purpose} \\
\midrule
src/App.jsx    & Main component -- form state, list render, delete handler \\
src/api.js     & \texttt{getStudents()}, \texttt{createStudent()}, \texttt{deleteStudent()} wrappers \\
src/index.css  & Full custom design system (CSS variables, grid layout, cards) \\
src/main.jsx   & React root entry point \\
vite.config.js & Vite build configuration with React plugin \\
\bottomrule
\end{tabularx}

\subsection{Frontend Features}

\begin{itemize}[noitemsep]
  \item \textbf{Add Student} -- form with Name, Email, Age, Course fields; submits POST
  \item \textbf{List Students} -- shows count badge, renders all enrolled students as cards
  \item \textbf{Delete Student} -- confirm dialog $\to$ DELETE request $\to$ removes from UI
  \item \textbf{Loading state} -- spinner while fetching initial data
  \item \textbf{Empty state} -- friendly message when no students are enrolled
  \item \textbf{Responsive layout} -- 2-column grid on desktop, single column on mobile ($<$768px)
\end{itemize}

\subsection{Running Frontend \& Backend Together}

\begin{lstlisting}[style=shellStyle, numbers=none, caption={Full stack development setup}]
# Terminal 1 - Start backend
cd "C:\C Projects\SEPQM\se3010-testing"
mvn spring-boot:run
# -> Backend running at http://localhost:8080

# Terminal 2 - Start frontend
cd "C:\C Projects\SEPQM\frontend"
npm install
npm run dev
# -> Frontend running at http://localhost:5173
# Open http://localhost:5173 in your browser
\end{lstlisting}

% ═══════════════════════════════════════════════════════════════════════════
\section{POM Dependency Configuration}
% ═══════════════════════════════════════════════════════════════════════════

\begin{lstlisting}[style=xmlStyle, caption={Key test dependencies in pom.xml}]
<!-- Spring Boot Test (JUnit FULLY EXCLUDED) -->
<dependency>
  <groupId>org.springframework.boot</groupId>
  <artifactId>spring-boot-starter-test</artifactId>
  <scope>test</scope>
  <exclusions>
    <exclusion>
      <groupId>org.junit.jupiter</groupId>
      <artifactId>junit-jupiter</artifactId>
    </exclusion>
    <exclusion>
      <groupId>org.junit.vintage</groupId>
      <artifactId>junit-vintage-engine</artifactId>
    </exclusion>
  </exclusions>
</dependency>

<!-- TestNG -->
<dependency>
  <groupId>org.testng</groupId>
  <artifactId>testng</artifactId>
  <version>7.10.2</version>
  <scope>test</scope>
</dependency>

<!-- Mockito standalone (no JUnit binding) -->
<dependency>
  <groupId>org.mockito</groupId>
  <artifactId>mockito-core</artifactId>
  <version>5.12.0</version>
  <scope>test</scope>
</dependency>

<!-- Cucumber BDD with TestNG runner -->
<dependency>
  <groupId>io.cucumber</groupId>
  <artifactId>cucumber-java</artifactId>
  <version>7.18.0</version>
  <scope>test</scope>
</dependency>
<dependency>
  <groupId>io.cucumber</groupId>
  <artifactId>cucumber-testng</artifactId>
  <version>7.18.0</version>
  <scope>test</scope>
</dependency>
<dependency>
  <groupId>io.cucumber</groupId>
  <artifactId>cucumber-spring</artifactId>
  <version>7.18.0</version>
  <scope>test</scope>
</dependency>
\end{lstlisting}

% ═══════════════════════════════════════════════════════════════════════════
\section{Troubleshooting}
% ═══════════════════════════════════════════════════════════════════════════

\begin{tabularx}{\textwidth}{>{\bfseries}lX}
\toprule
Problem & Solution \\
\midrule
\texttt{BUILD FAILURE} -- Cannot resolve symbol \texttt{testng}
  & Run \texttt{mvn dependency:resolve} then reload Maven in IDE \\
\addlinespace
\texttt{Connection refused} on frontend
  & Ensure backend is running (\texttt{mvn spring-boot:run}) before opening the UI \\
\addlinespace
PostgreSQL connection error during \texttt{mvn test}
  & Tests use \texttt{@ActiveProfiles("test")} which activates H2 -- no PostgreSQL needed for tests \\
\addlinespace
\texttt{H2 table not found}
  & Verify \texttt{src/test/resources/application-test.properties} has \texttt{ddl-auto=create-drop} \\
\addlinespace
Cucumber scenarios not found
  & Ensure \texttt{CucumberOptions\#features} path is relative to project root,
    and \texttt{glue = \{"steps"\}} matches the \texttt{steps} package \\
\addlinespace
\texttt{org.junit.*} import errors
  & Remove any leftover JUnit imports; all \texttt{@Test} must be from \texttt{org.testng.annotations} \\
\addlinespace
Port 8080 in use
  & Kill the process: \texttt{netstat -aon | findstr 8080} then \texttt{taskkill /PID <pid> /F} \\
\bottomrule
\end{tabularx}

% ═══════════════════════════════════════════════════════════════════════════
\section{Conclusion}
% ═══════════════════════════════════════════════════════════════════════════

This project demonstrates a complete, layered software testing strategy built entirely
on \textbf{TestNG} and \textbf{Cucumber-TestNG} -- with zero JUnit dependency:

\begin{itemize}[noitemsep]
  \item \textbf{Member 1} validates REST HTTP semantics (201, 200, 404, 400) using MockMvc with
        \texttt{org.testng.Assert} for all assertions.
  \item \textbf{Member 2} demonstrates test isolation, lifecycle management (\texttt{@BeforeMethod}/
        \texttt{@AfterMethod}), and data-driven testing via \texttt{@DataProvider} with Mockito
        standalone.
  \item \textbf{Member 3} brings business-readable BDD coverage via Gherkin feature files backed
        by MockMvc step definitions -- testing 201, 409 (duplicate email), and 400 scenarios.
  \item \textbf{Member 4} wires everything together via \texttt{AbstractTestNGCucumberTests},
        \texttt{testng.xml} suite orchestration, and Maven Surefire HTML reporting.
\end{itemize}

The \textbf{13 tests run in under 60 seconds}, with full HTML reports generated at
\texttt{target/site/surefire-report.html} and Cucumber-specific output at
\texttt{target/cucumber-reports/cucumber.html}.

\end{document}
